% Luc Destombe --- licence CC BY-NC-SA 4.0
% https://creativecommons.org/licenses/by-nc-sa/4.0/deed.fr
% Fichier autonome : compiler avec pdflatex, deux fois (signets, nombre de pages).
\documentclass[a4paper,10pt]{article}
\makeatletter
\newif\iflycee@niveaux
\lycee@niveauxtrue
\newif\iflycee@palatino
\lycee@palatinotrue

\RequirePackage[utf8]{inputenc}
\RequirePackage[T1]{fontenc}
\RequirePackage[french]{babel}
\RequirePackage{etoolbox}

\newcommand{\ShorthandsOff}{\@ifundefined{shorthandoff}{}{\shorthandoff{;:!?}}}
\newcommand{\ShorthandsOn}{\@ifundefined{shorthandon}{}{\shorthandon{;:!?}}}
\ShorthandsOff
\AtBeginDocument{\ShorthandsOn}
\AtBeginEnvironment{tikzpicture}{\ShorthandsOff}

\RequirePackage{geometry}
\RequirePackage{fancyhdr}
\RequirePackage{lastpage}
\RequirePackage{multicol}
\RequirePackage{array}
\RequirePackage{enumitem}
\RequirePackage{xcolor}
\RequirePackage{needspace}

\RequirePackage{amsmath,amssymb}
\RequirePackage{mathpazo}
\DeclareMathSizes{10.95}{10.95}{8}{5.5}
\begingroup\catcode`\_=\active \catcode`\^=\active
\gdef_#1{\sb{\scriptscriptstyle #1}}
\gdef^#1{\sp{\scriptscriptstyle\lycee@moinsepais #1}}
\endgroup
\newcommand\lycee@moins{\mathbin{%
  \setbox\z@\hbox{$\m@th\scriptscriptstyle\lycee@moinsorig$}%
  \dimen@=.55\wd\z@ \advance\dimen@-.5pt
  \hbox to.75\wd\z@{\hss$\m@th\scriptscriptstyle\vcenter{\hbox{%
    \tikz\draw[line width=.5pt,line cap=round](0,0)--(\the\dimen@,0);}}$\hss}}}
\begingroup\lccode`\~=`\- \lowercase{\endgroup\let~}\lycee@moins
\newcommand\lycee@moinsepais{\mathcode`\-="8000 }
\AtBeginDocument{\mathchardef\lycee@moinsorig=\mathcode`\-\relax}
\AtBeginDocument{\catcode`\_=12 \mathcode`\_="8000
  \catcode`\^=12 \mathcode`\^="8000 }
\newcommand\lycee@exposants{%
  \fontdimen13\textfont2=.48\fontdimen6\textfont2
  \fontdimen14\textfont2=.48\fontdimen6\textfont2
  \fontdimen15\textfont2=.48\fontdimen6\textfont2
  \fontdimen13\scriptfont2=.48\fontdimen6\scriptfont2
  \fontdimen14\scriptfont2=.48\fontdimen6\scriptfont2
  \fontdimen15\scriptfont2=.48\fontdimen6\scriptfont2}
\every@math@size\expandafter{\the\every@math@size\lycee@exposants}
\newlength{\retraitcalculs}
\setlength{\retraitcalculs}{2em}

\RequirePackage{tikz}
\usetikzlibrary{arrows.meta,calc,babel,shapes.misc}
\RequirePackage[most]{tcolorbox}
\tcbuselibrary{skins,breakable}

\definecolor{coulDef}{HTML}{1F4E79}
\definecolor{coulProp}{HTML}{7030A0}
\definecolor{coulRem}{HTML}{C55A11}
\definecolor{coulEx}{HTML}{2E7D32}
\definecolor{coulExo}{HTML}{B03A2E}
\definecolor{coulPart}{HTML}{1F4E79}
\definecolor{coulMeth}{HTML}{1B6E4F}
\definecolor{coulCourbe}{HTML}{0B5ED7}
\definecolor{coulCourbeB}{HTML}{C00000}
\definecolor{coulCourbeC}{HTML}{2E7D32}
\definecolor{coulGrille}{HTML}{8C8C8C}
\definecolor{coulRep}{HTML}{1B6E4F}
\definecolor{coulBar}{HTML}{2E7D32}

\newcommand{\R}{\mathbb{R}}
\newcommand{\Rs}{\mathbb{R}^{*}}
\renewcommand{\le}{\leqslant}
\renewcommand{\ge}{\geqslant}
\newcommand{\vect}[1]{\overrightarrow{#1}}
\newcommand{\norme}[1]{\left\lVert #1 \right\rVert}
\newcommand{\intervalle}[2]{\left[#1\,;#2\right]}
\RequirePackage{cancel}
\renewcommand{\CancelColor}{\color{coulExo}}
\newcommand{\fois}[1]{\times{\color{coulExo}#1}}
\newcommand{\barre}[1]{\cancel{#1}}

\tikzset{grille/.style={coulGrille,line width=0.4pt}}
\newcommand{\quadrillage}[4]{\draw[grille,step=1] (#1,#2) grid (#3,#4);}
\tikzset{
  croix/.style={cross out,draw,line width=0.6pt,inner sep=0pt,minimum size=4pt},
  croixdroite/.style={croix,rotate=45,line width=0.8pt},
}
\newcommand{\pt}[3]{\path (#1) node[croix]{} node[#2,font=\small,inner sep=2.5pt]{$#3$};}

\newcommand{\lycee@repere}[4]{%
  \draw[grille,step=1] (#1,#2) grid (#3,#4);
  \draw[-{Stealth[length=2mm]},thick] (#1,0)--({#3+0.4},0) node[below left=-1pt]{$x$};
  \draw[-{Stealth[length=2mm]},thick] (0,#2)--(0,{#4+0.4}) node[below left=-1pt]{$y$};
  \node[below left,font=\scriptsize,inner sep=1.5pt] at (0,0) {$0$};}
\newcommand{\lycee@gradx}[1]{\foreach \k/\l in {#1}{%
  \draw (\k,0.07)--(\k,-0.07) node[below,font=\scriptsize,inner sep=1.5pt]{$\l$};}}
\newcommand{\lycee@grady}[1]{\foreach \k/\l in {#1}{%
  \draw (0.07,\k)--(-0.07,\k) node[left,font=\scriptsize,inner sep=1.5pt]{$\l$};}}
\newcommand{\lycee@intff}[2]{\left[#1\,;#2\right]}
\newcommand{\lycee@intfo}[2]{\left[#1\,;#2\right[}
\newcommand{\lycee@intof}[2]{\left]#1\,;#2\right]}
\newcommand{\lycee@intoo}[2]{\left]#1\,;#2\right[}
\newcommand{\lycee@ens}[1]{\left\{#1\right\}}
\AtBeginDocument{%
  \providecommand{\repere}{\lycee@repere}%
  \providecommand{\gradx}{\lycee@gradx}%
  \providecommand{\grady}{\lycee@grady}%
  \providecommand{\Cf}{\mathcal{C}\sb{\scriptscriptstyle f}}%
  \providecommand{\Cg}{\mathcal{C}\sb{\scriptscriptstyle g}}%
  \providecommand{\intff}{\lycee@intff}%
  \providecommand{\intfo}{\lycee@intfo}%
  \providecommand{\intof}{\lycee@intof}%
  \providecommand{\intoo}{\lycee@intoo}%
  \providecommand{\ens}{\lycee@ens}}

\newlist{questions}{enumerate}{3}
\setlist[questions,1]{label=\arabic*\textdegree),leftmargin=*,itemsep=2pt,topsep=3pt}
\setlist[questions,2]{label=\alph*),leftmargin=*,itemsep=1pt,topsep=2pt}
\setlist[questions,3]{label=\roman*),leftmargin=*,itemsep=1pt,topsep=1pt}
\setlist[itemize]{label=\textbullet,leftmargin=*,itemsep=2pt,topsep=3pt}

\newcommand{\besoinplace}[1]{\par\penalty\@M\begingroup
  \dimen@=#1\relax\dimen@ii=\pagegoal\advance\dimen@ii-\pagetotal
  \ifdim\dimen@>\dimen@ii\ifdim\dimen@ii>\z@\vfil\fi\break\fi\endgroup}

\newcommand{\lycee@pied}{\thepage/\pageref{LastPage}}
\newcommand{\entete}[2]{%
  \pagestyle{fancy}\fancyhf{}%
  \fancyhead[L]{\footnotesize\color{coulPart}\textbf{#1}}%
  \fancyhead[R]{\footnotesize\color{coulPart}\textbf{#2}}%
  \fancyfoot[C]{\footnotesize\lycee@pied}%
  \renewcommand{\headrulewidth}{0.4pt}%
  \renewcommand{\headrule}{\hbox to\headwidth{%
    \color{coulPart!40}\leaders\hrule height \headrulewidth\hfill}}}

\RequirePackage{ccicons}
\newcommand{\lycee@licence}{{\scriptsize\color{gray}%
  \copyright\ Luc Destombe --- \ccbyncsa\ CC BY-NC-SA 4.0}}
\newif\iflycee@licence \lycee@licencetrue
\newcommand{\sanslicence}{\lycee@licencefalse}
\AtBeginDocument{\iflycee@licence\AddToHookNext{shipout/foreground}{%
  \put(0,0){\rlap{\kern\dimexpr1in+\hoffset+\oddsidemargin+\textwidth\relax
    \llap{\raisebox{-\dimexpr1in+\voffset+\topmargin+\headheight+\headsep
      +\textheight+\footskip\relax}{\lycee@licence}}}}}\fi}

\newcommand{\formulesserrees}{%
  \setlength{\abovedisplayskip}{3pt}\setlength{\belowdisplayskip}{3pt}%
  \setlength{\abovedisplayshortskip}{2pt}\setlength{\belowdisplayshortskip}{3pt}}

\setlength{\parindent}{0pt}

\RequirePackage[implicit=false,hidelinks,pdfencoding=auto,psdextra,bookmarksopen,
  bookmarksopenlevel=1,pdfcreator={},pdfproducer={}]{hyperref}
\RequirePackage{bookmark}
\hypersetup{pdfauthor={Luc Destombe},pdfkeywords={CC BY-NC-SA 4.0}}
\newcounter{lycee@signet}
\newcommand{\signet}[3][]{\stepcounter{lycee@signet}%
  \edef\lycee@dest{\if\relax\detokenize{#1}\relax
    signet.\arabic{lycee@signet}\else#1\fi}%
  \hypertarget{\lycee@dest}{}\bookmark[level=#2,dest=\lycee@dest]{#3}}
\pdfstringdefDisableCommands{%
  \def\R{R}\def\vect#1{#1}\def\Cf{Cf}\def\Cg{Cg}\def\fois#1{×#1}%
  \def\barre#1{#1}\def\intervalle#1#2{[#1;#2]}\def\intff#1#2{[#1;#2]}%
  \def\textsuperscript#1{#1}\def\dfrac#1#2{#1/#2}\def\frac#1#2{#1/#2}%
  \def\vec#1{#1⃗}}%
\RequirePackage[official]{eurosym}

\geometry{top=0.8cm,bottom=1.3cm,left=1cm,right=1cm,footskip=0.6cm}

\pagestyle{fancy}
\fancyhf{}
\renewcommand{\headrulewidth}{0pt}
\fancyfoot[C]{\footnotesize Page \thepage{} / \pageref{LastPage}}

\newcommand{\titrefiche}[2]{%
  \begin{center}
    \Large\textbf{#1}\\[2pt]
    \large\textbf{#2}\\[2pt]
  \end{center}
  \vspace{0.10cm}}

\setlength{\columnseprule}{0.4pt}
\setlength{\columnsep}{15pt}
\renewcommand{\columnseprulecolor}{\color{black!45}}
\setlength{\multicolsep}{2pt}

\newcounter{partiefiche}
\renewcommand{\thepartiefiche}{\Roman{partiefiche}}
\newif\ifapresbandeau
\newcommand{\lycee@bandeau}[1]{%
  \par\penalty-600
  \vspace{0.08cm}%
  \noindent\colorbox{coulPart}{%
    \parbox{\dimexpr\linewidth-2\fboxsep\relax}{%
      \color{white}\bfseries\raggedright\strut#1\strut}}%
  \nopagebreak\par\nopagebreak
  \vspace{0.06cm}\nopagebreak
  \apresbandeautrue}
\newcommand{\partiefiche}[1]{%
  \refstepcounter{partiefiche}%
  \lycee@bandeau{\signet{1}{\thepartiefiche{} --- #1}\thepartiefiche{} --- #1}}
\newcommand{\partiefichesansnum}[1]{\lycee@bandeau{\signet{1}{#1}#1}}

\newcounter{exo}
\newlength{\espaceexo}\setlength{\espaceexo}{7pt}
\newlength{\espacetitre}\setlength{\espacetitre}{2pt}
\newcommand{\exo}[1][\relax]{%
  \par
  \ifapresbandeau\apresbandeaufalse\else\penalty-150\addvspace{\espaceexo}\fi
  \refstepcounter{exo}%
  \noindent\signet[exercice-\arabic{exo}]{2}{Exercice \arabic{exo}}%
  \textbf{\color{coulExo}\underline{Exercice \theexo}}%
  \ifx\relax#1\relax\else\textbf{ : #1}\fi
  \par\nopagebreak\vspace{\espacetitre}\nopagebreak
  \ignorespaces}

\setlist[enumerate]{nosep,leftmargin=*,itemsep=2pt,topsep=2pt}
\setlist[itemize]{nosep,leftmargin=*,topsep=2pt}
\newcommand{\choix}[3]{\par\hspace*{1em}%
  \makebox[0.3\linewidth][l]{a) #1}\makebox[0.3\linewidth][l]{b) #2}c) #3}
\newcommand{\deux}[2]{\par\hspace*{0.6em}%
  \makebox[0.48\linewidth][l]{#1}#2}
\newcommand{\trois}[3]{\par\hspace*{0.6em}%
  \makebox[0.32\linewidth][l]{#1}\makebox[0.32\linewidth][l]{#2}#3}

  \renewcommand{\exo}[1][\relax]{%
    \ifapresbandeau\apresbandeaufalse\else\par\penalty-150\fi
    \refstepcounter{exo}%
    \vspace{0.05cm}%
    \noindent\signet[exercice-\arabic{exo}]{2}{Exercice \arabic{exo}}%
    \textbf{\color{coulExo}\underline{Exercice \theexo}}%
    \ifx\relax#1\relax\else\textbf{ : #1}\fi%
    \par\nopagebreak
    \ignorespaces}
  \newcommand{\rappel}[1]{%
    \par\vspace{0.06cm}\nopagebreak
    \noindent\fcolorbox{black!35}{black!5}{%
      \parbox{\dimexpr\linewidth-2\fboxsep-2\fboxrule\relax}{%
        \small\textbf{Rappel.} #1}}%
    \par\vspace{0.06cm}\nopagebreak}
  \setlist[enumerate]{nosep,leftmargin=*}
  \setlist[itemize]{nosep,leftmargin=*}
  \setlength{\multicolsep}{3pt plus 1pt minus 1pt}
  \setlength{\premulticols}{10pt}
  \setlength{\postmulticols}{10pt}
  \newlist{expr}{enumerate}{1}
  \setlist[expr]{label={\Alph*\ $=$},nosep,leftmargin=*,itemsep=0pt}

\renewenvironment{center}{\par\addvspace{3pt}\centering}{\par\addvspace{3pt}}
\formulesserrees
\AtBeginDocument{\formulesserrees}
\clubpenalty=10000
\widowpenalty=10000
\displaywidowpenalty=10000
\brokenpenalty=10000
\setlength{\parskip}{0pt}

\RequirePackage{ragged2e}
\setlength{\RaggedRightRightskip}{0pt plus 1fil}
\AtBeginDocument{\RaggedRight\binoppenalty=10000 \relpenalty=10000 }
\g@addto@macro\@parboxrestore{\RaggedRight}
\makeatother
\geometry{top=0.8cm,bottom=1.3cm,left=1cm,right=1cm,footskip=0.7cm,headheight=12pt}
\usepackage{tkz-tab}

\begin{document}
\titrefiche{1\textsuperscript{re} spécialité --- Dérivation}{Fiche d'exercices du chapitre}

\begin{multicols}{2}

\partiefiche{Limite en zéro d'une fonction}

\exo[Conjecturer avec un tableau de valeurs]
On considère $f(x)=\dfrac{x^{2}-3x}{x}$, avec $x\ne 0$.
\begin{enumerate}
  \item Calculer $f(x)$ pour $x$ égal à $-0{,}1$, puis $-0{,}01$, puis $0{,}01$ et
        $0{,}1$.
  \item Quelle semble être la limite de $f(x)$ quand $x$ tend vers $0$ ?
  \item Simplifier $f(x)$ et retrouver cette limite.
\end{enumerate}

\exo[Calculer des limites en zéro]
Calculer la limite en $0$ de chaque fonction.
\begin{multicols}{2}
\begin{enumerate}[label=\alph*)]
  \item $f(x)=\dfrac{x^{2}+4x}{x}$
  \item $g(x)=\dfrac{2x^{2}-6x}{x}$
  \item $h(x)=\dfrac{4x^{3}+x}{x}$
  \item $k(x)=\dfrac{3x^{3}-x^{2}+8x}{x}$
\end{enumerate}
\end{multicols}

\exo[Avec une identité remarquable]
Calculer la limite en $0$ de chaque fonction.
\begin{multicols}{2}
\begin{enumerate}[label=\alph*)]
  \item $f(x)=\dfrac{(1+x)^{2}-1}{x}$
  \item $g(x)=\dfrac{(5+x)^{2}-25}{x}$
  \item $h(x)=\dfrac{(x-2)^{2}-4}{x}$
  \item $k(x)=\dfrac{(3-x)^{2}-9}{x}$
\end{enumerate}
\end{multicols}

\exo[Limites quand $h$ tend vers $0$]
Calculer la limite de chaque expression quand $h$ tend vers $0$.
\begin{multicols}{3}
\begin{enumerate}[label=\alph*)]
  \item $\dfrac{h^{2}+6h}{h}$
  \item $\dfrac{(2+h)^{2}-4}{h}$
  \item $\dfrac{3(1+h)^{2}-3}{h}$
\end{enumerate}
\end{multicols}

\partiefiche{Nombre dérivé}

\exo[Nombre dérivé par la définition]
Pour chaque fonction, calculer le taux de variation $T(h)$, puis le nombre dérivé
demandé.
\begin{enumerate}[label=\alph*)]
  \item $f(x)=x^{2}$ : calculer $f'(3)$.
  \item $g(x)=x^{2}+2x$ : calculer $g'(1)$.
  \item $k(x)=-x^{2}+5x$ : calculer $k'(1)$.
  \item $m(x)=2x^{2}-3$ : calculer $m'(-1)$.
\end{enumerate}

\exo[Le taux de variation est donné]
Dans chaque cas, $h\ne 0$. En déduire le nombre dérivé demandé.
\begin{enumerate}[label=\alph*)]
  \item Le taux de variation de $f$ entre $1$ et $1+h$ est $T(h)=3h+7$ :
        calculer $f'(1)$.
  \item Le taux de variation de $g$ entre $4$ et $4+h$ est $T(h)=h^{2}-h-2$ :
        calculer $g'(4)$.
  \item On a $k(-1+h)-k(-1)=5h-2h^{2}$ : calculer $k'(-1)$.
\end{enumerate}

\exo[Avec des fractions]
\begin{enumerate}
  \item $f(x)=\dfrac{1}{x}$. Montrer que $f(4+h)-f(4)=\dfrac{-h}{4(4+h)}$, puis
        calculer $f'(4)$.
  \item $g(x)=\dfrac{1}{x+1}$. Calculer $g(1+h)-g(1)$, puis $g'(1)$.
\end{enumerate}

\exo[Pour un nombre $a$ quelconque]
On considère $f(x)=x^{2}-4x$ et un nombre $a$.
\begin{enumerate}
  \item Montrer que le taux de variation de $f$ entre $a$ et $a+h$ est
        $T(h)=2a-4+h$.
  \item En déduire $f'(a)$ en fonction de $a$.
  \item Calculer $f'(0)$, $f'(2)$ et $f'(5)$.
\end{enumerate}

\partiefiche{Tangente à une courbe}

\exo[Lecture graphique]
On a tracé la courbe $\Cf$ d'une fonction $f$ définie sur $\intff{-3}{5}$ et trois de
ses tangentes, aux points $A$, $B$ et $C$.
\begin{center}
\begin{tikzpicture}[x=0.5cm,y=0.5cm]
  \repere{-4}{-4}{6}{4}
  \gradx{-3,-2,-1,1,2,3,4,5} \grady{-3,-2,-1,1,2,3}
  \draw[coulCourbe,line width=1.1pt]
    (-3,-2) .. controls (-2.333,1) and (-1.667,2) .. (-1,2)
            .. controls (-0.333,2) and (0.333,1) .. (1,0)
            .. controls (1.667,-1) and (2.333,-2) .. (3,-2)
            .. controls (3.667,-2) and (4.333,-1) .. (5,2);
  \draw[coulExo,line width=0.9pt] (-2.6,2) -- (0.6,2);
  \draw[coulExo,line width=0.9pt] (-0.6,2.4) -- (2.6,-2.4);
  \draw[coulExo,line width=0.9pt] (1.4,-2) -- (4.6,-2);
  \path[coulCourbe] (-1,2) node[croix]{} (1,0) node[croix]{} (3,-2) node[croix]{};
  \node[font=\footnotesize] at (-1.5,2.5) {$A$};
  \node[font=\footnotesize,right] at (1.1,0.3) {$B$};
  \node[font=\footnotesize] at (3.5,-2.5) {$C$};
  \node[coulCourbe,font=\footnotesize] at (5.5,1.2) {$\Cf$};
\end{tikzpicture}
\end{center}
\begin{enumerate}
  \item Lire $f(-1)$, $f(1)$ et $f(3)$.
  \item Lire $f'(-1)$, $f'(1)$ et $f'(3)$.
  \item Donner une équation de la tangente à $\Cf$ en $B$.
\end{enumerate}

\exo[Équation de la tangente]
Dans chaque cas, donner une équation de la tangente à $\Cf$ au point d'abscisse $a$.
\begin{enumerate}[label=\alph*)]
  \item $a=2$ ; $f(2)=3$ et $f'(2)=4$
  \item $a=-1$ ; $f(-1)=2$ et $f'(-1)=-3$
  \item $a=0$ ; $f(0)=5$ et $f'(0)=\frac{1}{2}$
  \item $a=4$ ; $f(4)=-1$ et $f'(4)=0$
\end{enumerate}

\exo[Nombre dérivé, puis tangente]
\begin{enumerate}
  \item $f(x)=x^{2}+x-2$. Calculer $f'(1)$ par la définition, puis une équation de la
        tangente à $\Cf$ au point d'abscisse $1$.
  \item $g(x)=\dfrac{1}{x}$. Calculer $g'(-1)$ par la définition, puis une équation de
        la tangente à $\Cg$ au point d'abscisse $-1$.
\end{enumerate}

\exo[Tangente passant par un point]
\begin{enumerate}
  \item La tangente à $\Cf$ au point $A(1\,;-2)$ passe par $B(3\,;4)$. Calculer
        $f'(1)$, puis donner une équation de cette tangente.
  \item La tangente à $\Cg$ au point $C(-2\,;1)$ passe par $D(0\,;-3)$. Calculer
        $g'(-2)$, puis donner une équation de cette tangente.
\end{enumerate}

\partiefiche{Approximation linéaire}

\exo[Valeurs approchées sans calculatrice]
On utilise $f(a+h)\approx f(a)+f'(a)\times h$, pour $h$ proche de $0$ ; chaque
nombre dérivé se calcule par le taux d'accroissement.
\begin{enumerate}
  \item Pour $6{,}02^{2}$, on prend $f(x)=x^{2}$, $a=6$ et $h=0{,}02$. Calculer $f(6)$
        et $f'(6)$, puis une valeur approchée de $6{,}02^{2}$.
  \item Même méthode pour $1{,}01^{3}$, avec $g(x)=x^{3}$, $a=1$ et $h=0{,}01$.
  \item Pour $\dfrac{1}{0{,}98}$, avec $k(x)=\dfrac{1}{x}$ : choisir $a$ et $h$,
        calculer $k'(a)$, puis donner une valeur approchée.
  \item Même question pour $\dfrac{1}{2{,}01}$.
  \item Comparer les quatre résultats avec la calculatrice.
\end{enumerate}

\exo[Augmentations et baisses successives]
Chaque nombre dérivé se calcule par le taux d'accroissement.
\begin{enumerate}
  \item Un abonnement augmente de $1{,}5$~\% par mois pendant $2$ mois. Avec
        $f(x)=x^{2}$, calculer $f'(1)$, puis donner une valeur approchée de
        $1{,}015^{2}$ et le pourcentage d'augmentation sur les $2$ mois.
  \item La population d'un village baisse de $2$~\% par an pendant $3$ ans. Avec
        $g(x)=x^{3}$, calculer $g'(1)$, puis donner une valeur approchée de
        $0{,}98^{3}$ et le pourcentage de baisse sur les $3$ ans.
\end{enumerate}

\exo[Courbe et tangente]
\begin{enumerate}
  \item On considère $f(x)=x^{2}+3x$. Calculer $f(2)$, puis $f'(2)$ par le taux
        d'accroissement.
  \item En déduire une valeur approchée de $f(2{,}1)$, puis calculer l'écart avec la
        valeur exacte.
  \item Même question pour $f(1{,}95)$.
\end{enumerate}

\partiefiche{Dérivées des fonctions polynômes}

\exo[Dériver un polynôme]
Calculer la dérivée de chaque fonction.
\begin{multicols}{2}
\begin{enumerate}[label=\alph*)]
  \item $f(x)=7$
  \item $g(x)=5x-3$
  \item $h(x)=x^{6}$
  \item $k(x)=3x^{2}-4x+1$
  \item $m(x)=-2x^{3}+x^{2}-5x$
  \item $p(x)=x^{4}-2x^{3}+7x-1$
\end{enumerate}
\end{multicols}

\exo[Avec des fractions]
Calculer la dérivée de chaque fonction.
\begin{multicols}{2}
\begin{enumerate}[label=\alph*)]
  \item $f(x)=\frac{1}{2}x^{2}-3x$
  \item $g(x)=\frac{1}{3}x^{3}-x^{2}+4$
  \item $h(x)=\frac{1}{4}x^{4}+\frac{2}{3}x^{3}$
  \item $k(x)=\dfrac{x^{3}-6x}{3}$
\end{enumerate}
\end{multicols}

\exo[Développer d'abord]
Développer, puis calculer la dérivée de chaque fonction.
\begin{multicols}{2}
\begin{enumerate}[label=\alph*)]
  \item $f(x)=(x+4)^{2}$
  \item $g(x)=(3x-1)^{2}$
  \item $h(x)=x(x^{2}-5)$
  \item $k(x)=(x+2)(x-2)$
\end{enumerate}
\end{multicols}

\exo[Entraînement : polynômes]
Calculer la dérivée de chaque fonction.
\begin{multicols}{2}
\begin{enumerate}[label=\alph*)]
  \item $f(x)=x^{3}-7x^{2}+2x-9$
  \item $g(x)=-4x^{5}+3x^{2}$
  \item $h(x)=6x^{4}-x^{3}+x$
  \item $k(x)=\frac{2}{5}x^{5}-x$
  \item $m(x)=0{,}5x^{4}-1{,}5x^{2}$
  \item $p(x)=\dfrac{x^{2}-3}{4}$
  \item $q(x)=x^{7}-x^{5}+x^{3}$
  \item $r(x)=3(x^{2}-2x+5)$
\end{enumerate}
\end{multicols}

\exo[Fonction dérivée et nombres dérivés]
On considère $f(x)=x^{3}-3x^{2}+2$.
\begin{enumerate}
  \item Calculer $f'(x)$.
  \item En déduire $f'(0)$, $f'(1)$, $f'(-1)$ et $f'(2)$.
  \item Donner une équation de la tangente à $\Cf$ au point d'abscisse $1$.
  \item En quels points de $\Cf$ la tangente est-elle horizontale ?
\end{enumerate}

\partiefiche{Dérivées des autres fonctions de référence}

\exo[Nombres dérivés des fonctions de référence]
\begin{enumerate}[label=\alph*)]
  \item $f(x)=\dfrac{1}{x}$ : calculer $f'(5)$ et $f'\bigl(\frac{1}{2}\bigr)$.
  \item $g(x)=\sqrt{x}$ : calculer $g'(16)$ et $g'\bigl(\frac{1}{4}\bigr)$.
  \item $k(x)=\dfrac{1}{x^{2}}$ : calculer $k'(1)$ et $k'(-2)$.
\end{enumerate}

\exo[Dériver avec les fonctions de référence]
Calculer la dérivée de chaque fonction.
\begin{multicols}{2}
\begin{enumerate}[label=\alph*)]
  \item $f(x)=\dfrac{4}{x}$
  \item $g(x)=-3\sqrt{x}$
  \item $h(x)=\dfrac{1}{x^{4}}$
  \item $k(x)=\dfrac{5}{x^{3}}$
  \item $m(x)=2\cos x$
  \item $p(x)=-\sin x$
\end{enumerate}
\end{multicols}

\exo[Entraînement : fonctions de référence]
Calculer la dérivée de chaque fonction.
\begin{multicols}{2}
\begin{enumerate}[label=\alph*)]
  \item $f(x)=-\dfrac{7}{x}$
  \item $g(x)=8\sqrt{x}$
  \item $h(x)=\dfrac{2}{x^{2}}$
  \item $k(x)=-\dfrac{1}{x^{5}}$
  \item $m(x)=\dfrac{1}{3}\sqrt{x}$
  \item $p(x)=-4\cos x$
  \item $q(x)=\dfrac{1}{2x}$
  \item $r(x)=\dfrac{1}{2}\sin x$
\end{enumerate}
\end{multicols}

\exo[Tangentes aux courbes de référence]
\begin{enumerate}
  \item Donner une équation de la tangente à la courbe de $f(x)=\dfrac{1}{x}$ au point
        d'abscisse $1$.
  \item Donner une équation de la tangente à la courbe de $g(x)=\sqrt{x}$ au point
        d'abscisse $4$.
  \item Justifier que la courbe de $f$ n'a aucune tangente horizontale.
\end{enumerate}

\partiefiche{Dérivée d'une somme et d'un produit}

\exo[Dériver une somme]
Calculer la dérivée de chaque fonction.
\begin{multicols}{2}
\begin{enumerate}[label=\alph*)]
  \item $f(x)=x^{3}+\dfrac{1}{x}$
  \item $g(x)=2x^{2}-\sqrt{x}$
  \item $h(x)=\dfrac{3}{x}+5x$
  \item $k(x)=x+\cos x$
\end{enumerate}
\end{multicols}

\exo[Dériver un produit]
Calculer la dérivée de chaque fonction.
\begin{multicols}{2}
\begin{enumerate}[label=\alph*)]
  \item $f(x)=(3x-2)(x+4)$
  \item $g(x)=(x^{2}+1)(2x-3)$
  \item $h(x)=(x+1)\sqrt{x}$
  \item $k(x)=x^{2}\cos x$
\end{enumerate}
\end{multicols}

\exo[Entraînement : sommes et produits]
Calculer la dérivée de chaque fonction.
\begin{multicols}{2}
\begin{enumerate}[label=\alph*)]
  \item $f(x)=x^{2}-3x+\dfrac{4}{x}$
  \item $g(x)=4\sqrt{x}-2x^{3}$
  \item $h(x)=(x^{2}-4)(3x+1)$
  \item $k(x)=(2x-1)\sqrt{x}$
  \item $m(x)=(x+2)\cos x$
  \item $p(x)=\dfrac{1}{x}\,(x+5)$
  \item $q(x)=(x^{2}+x)(x^{2}-x)$
  \item $r(x)=\sin x\cos x$
\end{enumerate}
\end{multicols}

\exo[Deux façons de dériver]
On considère $f(x)=(2x+5)(1-3x)$.
\begin{enumerate}
  \item Calculer $f'(x)$ avec la dérivée d'un produit.
  \item Développer $f(x)$, puis dériver l'expression obtenue.
  \item Vérifier que les deux résultats sont égaux.
\end{enumerate}

\exo[Produit et tangente]
On considère $f(x)=(x-3)\sqrt{x}$ sur $\intoo{0}{+\infty}$.
\begin{enumerate}
  \item Montrer que $f'(x)=\dfrac{3x-3}{2\sqrt{x}}$.
  \item Calculer $f(4)$ et $f'(4)$.
  \item Donner une équation de la tangente à $\Cf$ au point d'abscisse $4$.
\end{enumerate}

\partiefiche{Dérivée d'un inverse et d'un quotient}

\exo[Dériver un inverse]
Calculer la dérivée de chaque fonction, là où elle est définie.
\begin{multicols}{2}
\begin{enumerate}[label=\alph*)]
  \item $f(x)=\dfrac{1}{2x+3}$
  \item $g(x)=\dfrac{1}{x^{2}+4}$
  \item $h(x)=\dfrac{3}{x-1}$
  \item $k(x)=\dfrac{1}{x^{2}-2x}$
\end{enumerate}
\end{multicols}

\exo[Dériver un quotient]
Calculer la dérivée de chaque fonction, là où elle est définie.
\begin{multicols}{2}
\begin{enumerate}[label=\alph*)]
  \item $f(x)=\dfrac{x+3}{x-2}$
  \item $g(x)=\dfrac{4x-1}{2x+1}$
  \item $h(x)=\dfrac{x^{2}}{x-3}$
  \item $k(x)=\dfrac{x^{2}-1}{x^{2}+1}$
\end{enumerate}
\end{multicols}

\exo[Entraînement : inverses et quotients]
Calculer la dérivée de chaque fonction, là où elle est définie.
\begin{multicols}{2}
\begin{enumerate}[label=\alph*),itemsep=4pt]
  \item $f(x)=\dfrac{1}{3-x}$
  \item $g(x)=\dfrac{-2}{x^{2}+x+1}$
  \item $h(x)=\dfrac{5x+2}{x-1}$
  \item $k(x)=\dfrac{1-2x}{x+3}$
  \item $m(x)=\dfrac{x^{2}+1}{x}$
  \item $p(x)=\dfrac{3x}{x^{2}+1}$
  \item $q(x)=\dfrac{x^{2}-4x}{x+2}$
  \item $r(x)=\dfrac{\sqrt{x}}{x+1}$
\end{enumerate}
\end{multicols}

\exo[Choisir la bonne formule]
Pour chaque fonction, dire s'il s'agit d'une somme, d'un produit, d'un inverse ou d'un
quotient, puis calculer sa dérivée.
\begin{multicols}{2}
\begin{enumerate}[label=\alph*)]
  \item $f(x)=3x^{4}-\dfrac{2}{x}$
  \item $g(x)=(x-1)(x^{2}+2)$
  \item $h(x)=\dfrac{4}{x^{2}+3}$
  \item $k(x)=\dfrac{x-4}{2x+1}$
  \item $m(x)=x^{2}\sqrt{x}$
  \item $p(x)=\sqrt{x}+\sin x$
\end{enumerate}
\end{multicols}

\exo[Tangentes à une hyperbole]
On considère $f(x)=\dfrac{2x-4}{x+1}$, pour $x\ne -1$.
\begin{enumerate}
  \item Montrer que $f'(x)=\dfrac{6}{(x+1)^{2}}$.
  \item $\Cf$ coupe l'axe des abscisses en $A$. Calculer les coordonnées de $A$, puis
        une équation de la tangente à $\Cf$ en $A$.
  \item Même question avec le point $B$ où $\Cf$ coupe l'axe des ordonnées.
\end{enumerate}

\exo[Deux fonctions, une même dérivée]
On considère $f(x)=\dfrac{2x+1}{x-1}$ et $g(x)=\dfrac{3}{x-1}$, pour $x\ne 1$.
\begin{enumerate}
  \item Calculer $f'(x)$ et $g'(x)$. Que remarque-t-on ?
  \item Calculer $f(x)-g(x)$. Expliquer la remarque de la question 1.
\end{enumerate}

\partiefiche{Dérivée des fonctions composées}

\exo[Dériver une puissance, une racine]
Calculer la dérivée de chaque fonction, là où elle est définie.
\begin{multicols}{2}
\begin{enumerate}[label=\alph*),itemsep=4pt]
  \item $f(x)=(3x+1)^{2}$
  \item $g(x)=(1-2x)^{2}$
  \item $h(x)=(x^{2}+3)^{3}$
  \item $k(x)=(4x-1)^{5}$
  \item $m(x)=\sqrt{2x+6}$
  \item $p(x)=\sqrt{x^{2}+4}$
\end{enumerate}
\end{multicols}

\exo[Cosinus, sinus et formule générale]
\begin{enumerate}
  \item Calculer la dérivée de chaque fonction.
\begin{multicols}{2}
\begin{enumerate}[label=\alph*)]
  \item $f(x)=\sin(3x)$
  \item $g(x)=\cos(5x+1)$
  \item $h(x)=\sin(x^{2})$
  \item $k(x)=\cos(1-x)$
\end{enumerate}
\end{multicols}
  \item On admet qu'une fonction $g$ a pour dérivée $g'(x)=\dfrac{1}{1+x^{2}}$. Avec
        la formule $f'(x)=u'(x)\times g'\big(u(x)\big)$, calculer la dérivée de
        $f(x)=g(2x)$.
\end{enumerate}

\exo[Composée et tangente]
On considère $f(x)=\sqrt{x^{2}+9}$, définie sur $\R$.
\begin{enumerate}
  \item Calculer $f'(x)$.
  \item Calculer $f(4)$ et $f'(4)$.
  \item Donner une équation de la tangente à $\Cf$ au point d'abscisse $4$.
  \item En déduire une valeur approchée de $f(4{,}1)$, puis la comparer avec la
        calculatrice.
\end{enumerate}

\partiefiche{Étude des variations d'une fonction}

\exo[Du signe de $f'$ aux variations de $f$]
Une fonction $f$ est dérivable sur $\intff{-4}{5}$ ; voici le signe de $f'(x)$.
\begin{center}
\begin{tikzpicture}
\tkzTabInit[lgt=1.6,espcl=1.5]{$x$/0.6,$f'(x)$/0.6}{$-4$,$-1$,$3$,$5$}
\tkzTabLine{,-,z,+,z,-,}
\end{tikzpicture}
\end{center}
On sait que $f(-4)=2$, $f(-1)=-3$, $f(3)=4$ et $f(5)=1$.
\begin{enumerate}
  \item Dresser le tableau de variations de $f$.
  \item Comparer $f(0)$ et $f(2)$.
\end{enumerate}

\exo[Variations de polynômes]
Dresser le tableau de variations de chaque fonction sur $\R$.
\begin{enumerate}[label=\alph*)]
  \item $f(x)=x^{2}-4x+7$
  \item $g(x)=-3x^{2}+12x-5$
  \item $h(x)=2x^{3}-3x^{2}-12x+1$
  \item $k(x)=x^{3}+3x-2$
\end{enumerate}

\exo[Variations d'un quotient]
On considère $f(x)=\dfrac{x^{2}+3}{x+1}$ sur $\intoo{-1}{+\infty}$.
\begin{enumerate}
  \item Montrer que $f'(x)=\dfrac{x^{2}+2x-3}{(x+1)^{2}}$.
  \item Étudier le signe de $f'(x)$ sur $\intoo{-1}{+\infty}$.
  \item Dresser le tableau de variations de $f$.
\end{enumerate}

\exo[Lire le signe de la dérivée]
On a tracé la courbe $\Cg$ d'une fonction $g$ dérivable sur $\intff{-4}{4}$. Ses
tangentes sont horizontales aux points d'abscisses $-2$, $1$ et $3$.
\begin{center}
\begin{tikzpicture}[x=0.5cm,y=0.5cm]
  \repere{-5}{-2}{5}{4}
  \gradx{-4,-3,-2,-1,1,2,3,4} \grady{-1,1,2,3}
  \draw[coulCourbe,line width=1.1pt]
    (-4,2) .. controls (-3.333,0.667) and (-2.667,-1) .. (-2,-1)
           .. controls (-1,-1) and (0,3) .. (1,3)
           .. controls (1.667,3) and (2.333,1) .. (3,1)
           .. controls (3.333,1) and (3.667,1.333) .. (4,2);
  \path[coulCourbe] (-2,-1) node[croix]{} (1,3) node[croix]{} (3,1) node[croix]{};
  \node[coulCourbe,font=\footnotesize] at (-4.4,2.8) {$\Cg$};
\end{tikzpicture}
\end{center}
\begin{enumerate}
  \item Dresser le tableau de signes de $g'(x)$ sur $\intff{-4}{4}$.
  \item Dresser le tableau de variations de $g$.
\end{enumerate}

\partiefiche{Extremums d'une fonction}

\exo[Lire les extremums]
Voici le tableau de variations d'une fonction $f$ dérivable sur $\intff{-2}{6}$.
\begin{center}
\begin{tikzpicture}
\tkzTabInit[lgt=1.4,espcl=1.5]{$x$/0.6,$f$/1.4}{$-2$,$1$,$4$,$6$}
\tkzTabVar{-/$1$,+/$5$,-/$-3$,+/$0$}
\end{tikzpicture}
\end{center}
\begin{enumerate}
  \item Donner le maximum et le minimum de $f$ sur $\intff{-2}{6}$.
  \item Donner le signe de $f'(x)$ selon les valeurs de $x$.
  \item Que vaut $f'(1)$ ? Et $f'(4)$ ?
\end{enumerate}

\exo[Extremums d'un polynôme]
On considère $f(x)=x^{3}-6x^{2}+9x+1$ sur $\intff{0}{5}$.
\begin{enumerate}
  \item Dresser le tableau de variations de $f$.
  \item Donner le maximum et le minimum de $f$ sur $\intff{0}{5}$.
  \item $f$ a-t-elle un extremum local en $1$ ? Lequel ?
\end{enumerate}

\exo[Vrai ou faux]
Dire si chaque affirmation est vraie ou fausse, et justifier.
\begin{enumerate}
  \item Si $f'(2)=0$, alors $f$ a un extremum en $2$.
  \item La fonction $g(x)=x^{3}+x$ n'a aucun extremum sur $\R$.
  \item La fonction $h(x)=x+\dfrac{1}{x}$ a pour minimum $2$ sur $\intoo{0}{+\infty}$.
\end{enumerate}

\partiefiche{Étude complète d'une fonction}

\exo[Étude d'un polynôme]
On considère $f(x)=-x^{3}+3x^{2}-1$ sur $\intff{-1}{3}$.
\begin{enumerate}
  \item Justifier que $f$ est définie et dérivable sur $\intff{-1}{3}$.
  \item Calculer $f'(x)$ et l'écrire sous forme factorisée.
  \item Étudier le signe de $f'(x)$, puis dresser le tableau de variations de $f$.
  \item Donner les extremums locaux de $f$, puis son maximum et son minimum sur
        $\intff{-1}{3}$.
  \item Donner une équation de la tangente $T$ à $\Cf$ au point d'abscisse $1$.
  \item Avec un tableau de valeurs (pas de $0{,}5$), tracer $T$ et $\Cf$ dans un repère
        orthonormé d'unité $1$~cm.
\end{enumerate}

\exo[Étude d'un quotient défini sur $\R$]
On considère $f(x)=\dfrac{2x+1}{x^{2}+2}$.
\begin{enumerate}
  \item Justifier que $f$ est définie sur $\R$.
  \item Montrer que $f'(x)=\dfrac{-2(x+2)(x-1)}{(x^{2}+2)^{2}}$.
  \item Étudier le signe de $f'(x)$, puis dresser le tableau de variations de $f$.
  \item Donner les extremums locaux de $f$.
  \item Donner une équation de la tangente $T$ à $\Cf$ au point d'abscisse $0$.
  \item Avec un tableau de valeurs (pas de $1$), tracer $T$ et $\Cf$ sur
        $\intff{-4}{4}$, dans un repère orthonormé d'unité $2$~cm.
\end{enumerate}

\exo[Étude d'un quotient avec une valeur interdite]
On considère $f(x)=\dfrac{x^{2}-3}{x-2}$.
\begin{enumerate}
  \item Déterminer l'ensemble de définition de $f$.
  \item Montrer que $f'(x)=\dfrac{(x-1)(x-3)}{(x-2)^{2}}$.
  \item Étudier le signe de $f'(x)$, puis dresser le tableau de variations de $f$.
  \item Donner les extremums locaux de $f$. Que remarque-t-on ?
  \item Donner une équation de la tangente $T$ à $\Cf$ au point d'abscisse $0$.
  \item Avec un tableau de valeurs, tracer $T$ et $\Cf$ sur $\intff{-3}{1{,}5}$ et sur
        $\intff{2{,}5}{7}$, dans un repère orthonormé d'unité $1$~cm.
\end{enumerate}

\columnbreak
\partiefiche{Problèmes concrets}

\exo[Le vol d'une montgolfière]
Une montgolfière décolle d'un village de montagne à 10~h~50. Son altitude, en centaines
de mètres, est $f(t)=\dfrac{8t+6}{t^{2}+1}$, où $t\in\intff{0}{5}$ est le temps en
heures écoulé depuis le décollage.
\begin{enumerate}
  \item À quelle altitude se trouve le village ?
  \item Montrer que $f'(t)=\dfrac{-4(2t-1)(t+2)}{(t^{2}+1)^{2}}$.
  \item Étudier le signe de $f'(t)$ sur $\intff{0}{5}$, puis dresser le tableau de
        variations de $f$.
  \item À quelle heure la montgolfière commence-t-elle à perdre de l'altitude ? Quelle
        est alors son altitude ?
  \item À quelle heure revient-elle à l'altitude du village ?
\end{enumerate}

\exo[Des bactéries et un antibiotique]
Un antibiotique est ajouté à un milieu de culture. Le nombre de bactéries, en milliers,
est $N(t)=t^{3}-12t^{2}+36t+10$, où $t\in\intff{0}{8}$ est le temps en heures. La
vitesse de croissance de la population, en milliers de bactéries par heure, est
$N'(t)$.
\begin{enumerate}
  \item Calculer $N(0)$ et interpréter le résultat.
  \item Montrer que $N'(t)=3(t-2)(t-6)$.
  \item Dresser le tableau de variations de $N$ sur $\intff{0}{8}$.
  \item Calculer $N'(0)$ et $N'(4)$, puis interpréter ces deux nombres.
  \item Montrer que $N(t)-10=t(t-6)^{2}$. Le nombre de bactéries peut-il descendre
        sous sa valeur de départ ?
\end{enumerate}

\exo[Le coût moyen d'un artisan]
Un artisan fabrique entre $5$ et $50$ lampes par mois. Le coût moyen d'une lampe, en
euros, est $M(x)=x+\dfrac{400}{x}$, où $x\in\intff{5}{50}$ est le nombre de lampes.
\begin{enumerate}
  \item Calculer $M(10)$ et interpréter le résultat.
  \item Montrer que $M'(x)=\dfrac{(x-20)(x+20)}{x^{2}}$.
  \item Dresser le tableau de variations de $M$. Combien de lampes faut-il fabriquer
        pour que le coût moyen soit le plus bas ? Quel est-il ?
  \item Vrai ou faux ? Justifier. \og Le coût moyen est d'au plus $50$~\euro{} dès
        que l'artisan fabrique entre $10$ et $40$ lampes. \fg
  \item Vrai ou faux ? Justifier. \og La tangente à la courbe de $M$ au point
        d'abscisse $10$ est parallèle à la droite d'équation $y=-2x+1$. \fg
\end{enumerate}

\end{multicols}
\end{document}
